\section{Data} \label{sec:data}

To evaluate the escape velocity of the Milky Way, we use the third catalog from the \Gaia space mission (\Gaia~DR3) \citep{gaia:2016,gaiaDR3}, which contains $\sim 33$ million stars with 6D kinematics, an increase of about a factor of five from the previous release of \Gaia~eDR3 \citep{GaiaeDR3}. 
As we are interested in modeling the tail of the speed distribution, we extract a high-quality sample of stars from this catalog in order to avoid stars with contaminated radial velocities and/or unphysical parallax measurements. Such contaminants could bias escape velocity estimates toward higher values since these analyses are sensitive to the fastest star in the sample (see e.g. \cite{Koppelman2021}). 

To obtain a high quality sample, we follow the selection criteria of \cite{marchetti}, imposing quality cuts on parallax measurements, renormalized unit weight error (RUWE), number of transits used to calculate the radial velocity of the source, and signal-to-noise ratio. We do however demand  higher-quality signal to noise ratios of $10$ (as opposed to $5$) in line with \cite{gaiaRVs}, which found a large number of spurious radial velocity measurements near the Solar position at low signal-to-noise, and parallaxes with at most $10\%$ error (as opposed to $20\%$) to reduce the contamination of stars across radial bins. 
For parallax measurements, we apply a zero-point correction according to \cite{zpcorr} and choose not to include those stars for which a zero-point-corrected and positive parallax is unavailable. Note that an alternative implementation of the cut on parallax percentage errors would be a ``bin membership" cut whereby a star only be included in the analysis if it is a member of its hosting radial bin at some chosen confidence level. In the present analysis, such a cut would be too restrictive and reduce statistics away from the solar position due to the limitations of parallax measurements. This may lead to some leakage of stars near bin edges across radial bins, but since the escape velocity is not expected to vary significantly over distances $\lesssim 0.5\,\rm{kpc}$, we do not expect this to appreciably bias results as is shown in Fig. \ref{fig:offsetbins}, in which we repeated our analysis with a shift in the binning. 

We also require that the derived percentage error on the speed be less than $1\%$ since large errors can significantly modify the shape of the tail of the distribution. In some works such as \cite{Piffl2014, Monari2018}, only the retrograde stars in the sample are modeled to avoid disk contamination, but following \cite{Necib2022} we can also remove disk stars via high minimum speed cuts, here not considering stars with speeds below $300\,\rm{km\,s^{-1}}$, or by choosing more flexible models that can account for the presence of additional velocity components.

The cuts are summarized as follows:
\setlist{nolistsep}
\begin{itemize}[itemsep=0em, label=-]
    \item \texttt{RUWE} $\leq 1.4$
    \item \texttt{RV\_NB\_TRANSITS} $\geq10$   
    \item \texttt{RV\_EXPECTED\_SIG\_TO\_NOISE} $\geq10$   
    \item Positive zero-point-corrected parallax $\varpi - \varpi_{zp} > 0$
    \item Parallax error: $\sigma_{\varpi}/(\varpi - \varpi_{zp}) \leq 10\%$
    \item Minimum speed: $v\geq 300\,\rm{km\,s^{-1}}$ 
    \item Speed error: $\sigma_v/v \leq 1\%$
\end{itemize}

The data processing pipeline\footnote{Pipeline and data products are available at \url{https://doi.org/10.5281/zenodo.8088365}.} used to calculate the stellar kinematics and associated uncertainties in Galactocentric coordinates is similar to that of \cite{Necib2022}, with the addition of the zero-point correction of parallaxes and stricter quality cuts. In particular, we use sources with parallax errors of $10\%$ or less (rather than $20\%$) to reflect the finer radial binning used in this analysis ($1\,\rm{kpc}$ bins as opposed to $2\,\rm{kpc}$), and percentage speed errors of at most $1\%$ (rather than $5\%$) because speed measurement errors can dominate the shape of the high-speed tail; these choices are enabled by the improved statistics of \Gaia~DR3. The Solar position used in this analysis is $x_\odot = -8.122\,\rm{kpc}$ \citep{gravity:2018}, $y_\odot = 0\,\rm{kpc}$, and $z_\odot=0.0208\,\rm{kpc}$ \citep{Bennett:2018} in Galactocentric coordinates. For the Solar peculiar velocity vector we use $v_\odot = (12.9, 245.6, 7.78)\,\rm{km\,s^{-1}}$ \citep{Drimmel:2018}. 

The cuts we utilize result in one of the highest-quality 6D kinematics datasets used for the escape velocity measurements to date. It is also the largest, containing $12,317$ stars above a speed of $300\,\rm{km\,s^{-1}}$, compared to the $3,932$ within $1\,\rm{kpc}$ of the solar position of \cite{Necib2022}, and the $2,067$ in the 6D kinematics sample of \cite{Koppelman2021}. We represent this selection of stars in \Gaia~DR3 in Figure \ref{fig:data}, along with one of the escape velocity profiles obtained in this paper, discussed in detail in Sec. \ref{sec:results}. 

In Figure \ref{fig:data}, one can note some general features of the data. First, there is a slight overdensity of high velocity stars near the solar position. This excess may be due to a larger number of spurious radial velocity measurements in this region \citep{gaiaRVs}, as this feature appears primarily close to the solar position. 
%

Second, the speed of the fastest stars generally decreases from $\sim 4.5$ to $12\,\rm{kpc}$, partly as a result of the increasing Galactic potential. At Galactocentric radii smaller than $4\,\rm{kpc}$, however, there is a decrease in statistics resulting from parallax quality cuts, leading to fewer high speed stars. In general, the number of observed stars in each radial bin is modulated by the \Gaia~selection function, which is typically modeled as a function of magnitude and sky position \citep{gaiaeDR3selection}, or a combined metric of both parameters such as the median magnitude of a patch of sky \citep{gaiaDR3selection}. 
Assuming no significant correlations between completeness and speed, we  expect the shape of the speed distribution in a given radial bin is not strongly biased by the selection function. However, even if the selection function impacts the detailed shape of the distribution, our goal is to extract the cutoff of the speed distribution and not its exact shape.

One can also note over and under-densities in the contours from $7-8\,\rm{kpc}$ at speeds close to and below $300\,\rm{km\,s^{-1}}$. The source of this feature is the resonant band structure of the speed distribution of the Milky Way \citep{Antoja2018}, and in fact there are other parallel bands which are simply not visible in this limited data sample. Over and under-densities at certain Galactocentric radii cannot be modeled by simple power law distributions, and thus fits to single power laws could lead to biased escape velocity estimates.
This already suggests the need for models that are robust to the presence of such features.
